\apartado{Arquitectura propuesta y flujo de información}
La figura representa el destino arquitectónico de la v1. Las fases H0 y H1 desarrollan la base y el espacio académico; los componentes de IA, cola y almacenamiento documental se muestran para explicar sus relaciones futuras, sin atribuirles implementación completa.

\begin{center}
\begin{tikzpicture}[>=Latex,node distance=0.75cm and 0.7cm,
 box/.style={draw=Azul,line width=0.9pt,fill=Marfil,align=center,minimum height=0.95cm,text width=3.25cm,font=\small},
 future/.style={box,dashed,fill=white},arr/.style={->,draw=Azul,thick}]
\node[box] (user) {Estudiante\\Navegador web};
\node[box,right=of user] (auth) {Firebase Auth\\Identidad de persona};
\node[box,below=of user] (web) {Web React\\Contenedor Cloud Run};
\node[box,below=of web] (api) {API FastAPI\\Actor y autorización};
\node[box,right=of api] (db) {Firestore\\Registros y metadatos};
\node[future,below=of api] (tasks) {Cloud Tasks\\Despacho\\autenticado};
\node[future,right=of tasks] (store) {\mbox{Cloud Storage}\\privado\\Originales y derivados};
\node[future,below=of tasks] (worker) {Worker HTTP privado\\Cloud Run};
\node[future,right=of worker] (ai) {Proveedor DeepSeek\\Adaptador del servidor};
\node[future,left=of worker] (renderer) {Renderer separado\\VM limitada + SDK};
\draw[arr] (user)--(web);
\draw[arr] (user)--(auth);
\draw[arr] (web)--node[left,font=\scriptsize]{token}(api);
\draw[arr] (api)--(db);
\draw[arr] (api)--(tasks);
\draw[arr] (tasks)--(worker);
\draw[arr] (api)--(store);
\draw[arr] (worker)--(store);
\draw[arr] (worker)--(ai);
\draw[arr] (worker.south) -- ++(0,-0.35) -| ([xshift=0.55cm]db.east) -- (db.east);
\draw[arr] (web.west) to[out=180,in=90] (renderer.north);
\end{tikzpicture}
\end{center}
\begin{center}\small\textbf{Figura 1.} Arquitectura de destino. Trazo discontinuo: capacidades de integración posterior. Elaboración propia a partir del plan del proyecto.\end{center}

\subsection{Recorrido síncrono y recorrido diferido}
En H1, el navegador obtiene una sesión y solicita datos al API. El servidor valida el token, habilitación del piloto y pertenencia del recurso. Consulta o modifica registros y devuelve un resultado con revisión y errores definidos. Los datos persistentes no deben depender de la memoria de una instancia de contenedor.

Posteriormente, una operación lenta creará un trabajo persistente. La cola notificará al worker, que verificará el alcance y actualizará el estado. El navegador consultará el progreso: cerrar una pestaña no debe perder el trabajo. Un recurso interactivo autorizado se enviará al renderer en otro origen; ese componente no recibirá secretos ni acceso general a documentos.

\textbf{Fronteras principales:} navegador--API (actor humano), API--datos (propiedad), Tasks--worker (identidad de servicio), servidor--IA (contenido mínimo autorizado) y aplicación--renderer (mensajes acotados). Cada frontera requiere pruebas específicas; una flecha en un diagrama no acredita el permiso.
